\section{Completed derivatives and their weights}
\label{sec:weights}

A large analytic rank forces a completed derivative to vanish. For the
matching functional-equation sign, we express that vanishing as an exact
weighted-series identity, partition the series into short pieces, and
bound the remaining infinite piece using the second moments of
Section~\ref{sec:analytic}.
The short pieces will be treated by mollification in
Section~\ref{sec:model}.

Fix a curve \(F\) in the finite family \(\mathcal F\), write
\(\lambda=\lambda_F\), and put \(C_F=\sqrt{q_F}/(2\pi)\).
The variable \(x\) will range over a fixed compact subset
\(\mathcal K\) of \(\mathbb R\setminus\{0\}\), chosen large enough to
contain the support of the averaging weight in
Proposition~\ref{analytic:comparison}. Derivative bounds in \(x\) are
understood on a fixed compact neighborhood of \(\mathcal K\) disjoint
from zero. All constants below may depend on this compact set and on
the fixed finite family.

The order of choices is as follows. First fix the exponents in
Proposition~\ref{analytic:polynomial-moment}, and all the finite orders of
differentiation needed for Proposition~\ref{analytic:comparison}, the
Mellin inversion below, and the subsequent Fourier estimates. Next
choose a sufficiently large fixed integer \(k_0\). Finally choose one
lower bound for \(X\), valid simultaneously for all the integers
\begin{equation}
 \ell=\log X,\qquad
 k_0\le k\le \lfloor\ell^{3/5}\rfloor .
 \label{weights:range}
\end{equation}
The estimates in this section respect this order. In particular, the
number of derivatives is never allowed to increase with \(k\) or \(X\).

\subsection{The exact completed derivative identity}

For \(k\ge1\), \(x\ne0\), and real \(u\), define
\begin{equation}
 W_{k,x}(u)=
 \int_0^\infty e^{-v}
 \left(1-u+\frac{\log(C_F|x|v)}{\ell}\right)_+^k\,dv,
 \qquad y_+=\max(y,0).
 \label{weights:W}
\end{equation}
For every fixed choice of the parameters the integral is finite.
The weighted sum below is an infinite, absolutely convergent smoothed
Dirichlet series; its individual short pieces will be finite sums.

\begin{lemma}[Completed derivative identity]
\label{weights:vanishing}
Let \(d\) be admissible, \(X\le |d|<2X\), and \(x=d/X\). Set
\[
 g_d(z)=L_{F^{(d)}}(1/2+z)(C_F|d|)^z\Gamma(1+z),
 \qquad \epsilon_d=\operatorname{sign}(F^{(d)}).
\]
Then \(g_d\) is entire, \(g_d(-z)=\epsilon_d g_d(z)\), and
\begin{equation}
 \begin{split}
 \mathcal S_d
 &:=
 \sum_{n\ge1}\frac{\lambda(n)\chi_d(n)}{\sqrt n}
                 W_{k,x}(\ell^{-1}\log n)\\
 &=\frac{k!}{\ell^k}\frac{1}{2\pi i}
       \int_{(2)}g_d(z)\frac{dz}{z^{k+1}}.
 \end{split}
 \label{weights:completed-integral}
\end{equation}
More precisely,
\begin{equation}
 \bigl(1+\epsilon_d(-1)^k\bigr)\mathcal S_d
       =\ell^{-k}g_d^{(k)}(0).
 \label{weights:exact-derivative}
\end{equation}
Consequently, if \(a(F^{(d)})>k\) and
\(\epsilon_d=(-1)^k\), then
\begin{equation}
 \sum_{n\ge1}\frac{\lambda(n)\chi_d(n)}{\sqrt n}
                  W_{k,d/X}(\ell^{-1}\log n)=0.
 \label{weights:zero}
\end{equation}
\end{lemma}

\begin{proof}
The completed function is
\[
 \Lambda_d(s)=(C_F|d|)^s\Gamma(s+1/2)L_{F^{(d)}}(s).
\]
Thus \(g_d(z)=(C_F|d|)^{-1/2}\Lambda_d(1/2+z)\).
Its entireness and reflection law follow from those of \(\Lambda_d\).

For \(T>0\), Mellin inversion gives
\[
 \frac{1}{2\pi i}\int_{(2)}\frac{T^z}{z^{k+1}}\,dz
       =\frac{(\log T)_+^k}{k!}.
\]
Insert the absolutely convergent Dirichlet series for
\(L_{F^{(d)}}(1/2+z)\), and use
\(\Gamma(1+z)=\int_0^\infty e^{-v}v^z\,dv\).
On this line all exchanges are absolutely convergent: the series is
dominated by \(\sum_n\tau(n)n^{-5/2}\), and the other two factors to be
integrated are dominated by
\(e^{-v}v^2\) and \(|2+it|^{-k-1}\).
The result is \eqref{weights:completed-integral}, because
\[
 \ell^{-1}\log(C_F|d|v/n)
       =1-\ell^{-1}\log n+\ell^{-1}\log(C_F|x|v).
\]
This argument, or the same computation with absolute coefficient
values, also proves absolute convergence of the series.

Let \(I_+\) and \(I_-\) denote the integrals of
\(g_d(z)z^{-k-1}/(2\pi i)\) on the upward oriented lines
\(\operatorname{Re}z=2\) and \(\operatorname{Re}z=-2\).
Shifting between these lines crosses only the pole at zero, and gives
\[
 I_+-I_-=\frac{g_d^{(k)}(0)}{k!}.
\]
The horizontal integrals tend to zero: in this fixed strip the
uncompleted \(L\)-function has polynomial growth, while the gamma
factor has exponential decay on horizontal segments tending to
infinity. Apparent gamma poles on the real axis are removable in
\(g_d\). Changing \(z\) to \(-z\), with orientations included, gives
\[
 I_-=(-1)^{k+1}\epsilon_d I_+.
\]
Multiplication by \(k!/\ell^k\) proves
\eqref{weights:exact-derivative}. Finally, the completing factor is
holomorphic and nonzero at zero. Hence \(a(F^{(d)})>k\) implies
\(g_d^{(k)}(0)=0\), and the stipulated sign makes the coefficient on
the left of \eqref{weights:exact-derivative} equal to two.
\end{proof}

\subsection{Partition and derivative estimates}

Fix a \(C^\infty\), nondecreasing function \(f_*\colon\mathbb R\to[0,1]\)
which is zero on \((-\infty,1/4]\) and one on \([1/2,\infty)\).
Define
\begin{equation}
 J=\left\lfloor\frac{\log\ell}{10\log2}\right\rfloor,\qquad
 \delta_j=2^{-j},\qquad
 F_j(u)=f_*((1-u)/\delta_j)\quad(0\le j\le J).
 \label{weights:scales}
\end{equation}
In particular, for sufficiently large \(X\),
\begin{equation}
 J\ge1,\qquad
 \ell^{-1/10}\le\delta_J<2\ell^{-1/10},\qquad
 \frac{k}{\ell\delta_j}\le\ell^{-3/10}
       \quad(0\le j\le J).
 \label{weights:scale-bounds}
\end{equation}
Set
\begin{equation}
 \begin{aligned}
 w_{0,x}(u)&=F_0(u)W_{k,x}(u),\\
 w_{j,x}(u)&=(F_j(u)-F_{j-1}(u))W_{k,x}(u)
                    &&(1\le j\le J),\\
 w_{*,x}(u)&=(1-F_J(u))W_{k,x}(u).
 \end{aligned}
 \label{weights:partition}
\end{equation}
These weights sum exactly to \(W_{k,x}\). For \(V=\chi_d,z_m\), or \(Z\),
write
\begin{equation}
 \begin{aligned}
 S_{j,V}(x)&=\sum_{n\ge1}
       \frac{\lambda(n)V(n)}{\sqrt n}w_{j,x}(\ell^{-1}\log n),\\
 S_{*,V}(x)&=\sum_{n\ge1}
       \frac{\lambda(n)V(n)}{\sqrt n}w_{*,x}(\ell^{-1}\log n).
 \end{aligned}
 \label{weights:S}
\end{equation}
The first definition includes \(j=0\). All the formal character
values in these definitions have modulus at most one. The sums
through \(j=J\) have length at most \(X^{1-\delta_j/4}\), and
\begin{equation}
 \operatorname{supp} w_{j,x}
       \subset[1-\delta_j,1-\delta_j/4]\quad(1\le j\le J).
 \label{weights:short-support}
\end{equation}
Indeed \(F_j-F_{j-1}\) vanishes when \(1-u\le\delta_j/4\) or
\(1-u\ge\delta_j\). The terminal weight is supported in
\(u\ge1-\delta_J/2\).

\begin{lemma}[Fixed-order weight bounds]
\label{weights:derivatives}
Fix an integer \(r\ge0\), and then take \(k_0\ge r+2\).
For \(a,b\ge0\), \(a+b\le r\), uniformly in \eqref{weights:range},
\(x\in\mathcal K\), and the indicated arguments, one has
\begin{align}
 |\partial_u^a\partial_x^b w_{j,x}(u)|
   &\le C_r(k/\delta_j)^r\delta_j^k
                  &&(1\le j\le J,\ u\in\mathbb R), \label{weights:short-derivatives}\\
 |\partial_u^a\partial_x^b w_{0,x}(u)|
   &\le C_r k^r
                  &&(u\ge-\log2/\ell). \label{weights:main-derivatives}
\end{align}
There is also \(c>0\), independent of \(r,k,X\), such that, with
\(\delta=\delta_J\) and \(n=e^{\ell u}>0\),
\begin{equation}
 |\partial_u^a\partial_x^b w_{*,x}(u)|
       \le C_r(k/\delta)^r\delta^k e^{-cn/X}.
 \label{weights:terminal-derivatives}
\end{equation}
Changing the fixed compact neighborhood of \(\mathcal K\) changes
only the constants.
\end{lemma}

\begin{proof}
It is convenient first to differentiate using
\(D_x=x\partial_x\). For
\[
 h=1-u+\ell^{-1}\log(C_F|x|v)
\]
and \(a+b\le r\), differentiation under the integral gives
\begin{equation}
 \partial_u^aD_x^b h_+^k
    =(-1)^a\ell^{-b}(k)_{a+b}h_+^{k-a-b},
 \label{weights:power-derivative}
\end{equation}
where \((k)_t=k(k-1)\cdots(k-t+1)\).
The choice \(k\ge r+2\) ensures enough continuous derivatives at
\(h=0\); domination for differentiating follows from the estimates
below. Ordinary \(x\)-derivatives are fixed linear combinations of
\(x^{-b}D_x^i\), \(i\le b\), and therefore satisfy the same bounds on
our compact set.

Suppose \(1-u\le\delta'\), where \(\delta'>0\). For \(0\le t\le r\),
write \(m=k-t>0\) and \(a_t=m/(\ell\delta')\). The elementary
inequality \((1+q)_+^m\le e^{mq}\), valid for every real \(q\), gives
\begin{equation}
 h_+^m
    \le(\delta')^m(C_F|x|v)^{a_t}.
 \label{weights:gamma-majorant}
\end{equation}
If \(\delta'=\delta_j\), then \(0<a_t\le\ell^{-3/10}\) by
\eqref{weights:scale-bounds}, so
\[
 \int_0^\infty e^{-v}(C_F|x|v)^{a_t}\,dv
       =(C_F|x|)^{a_t}\Gamma(1+a_t)\le C.
\]
The bound is uniform once \(X\) is large enough that \(a_t\le1\).
Thus, on \(1-u\le\delta_j\), a derivative of \(W_{k,x}\) involving at
most \(r\) differentiations is bounded by a constant times
\(k^r\delta_j^{k-r}\).
The cutoff derivatives of order \(q\) cost \(O_r(\delta_j^{-q})\).
More precisely, each Leibniz term with \(t\) derivatives hitting
\(W\) and \(q\) hitting a cutoff has \(t+q\le r\), and is at most
\[
 C_r k^t\delta_j^{k-t-q}
       \le C_r(k/\delta_j)^r\delta_j^k.
\]
This proves \eqref{weights:short-derivatives}.

For \(u\ge-\log2/\ell\), use instead
\(\delta'=1+\log2/\ell\). Now
\((\delta')^k\le\exp(k\log2/\ell)\le2\) for large \(X\), and
\(a_t\le k/\ell\le\ell^{-2/5}\).
The same argument and the fixed derivative bounds for \(F_0\)
prove \eqref{weights:main-derivatives}.

On the support of the terminal cutoff, or any nonzero derivative
of it, \(1-u\le\delta_J/2\le\delta_J\).
Moreover \(h>0\) implies
\[
 v>\frac{n}{C_F|x|X}.
\]
Retain half the exponential in the integral. Uniformly on our
compact set,
\begin{equation}
 \begin{split}
 \int_{\{h>0\}}e^{-v}(C_F|x|v)^{a_t}\,dv
 &\le e^{-n/(2C_F|x|X)}
       (C_F|x|)^{a_t}
       \int_0^\infty e^{-v/2}v^{a_t}\,dv\\
 &=e^{-n/(2C_F|x|X)}
       (C_F|x|)^{a_t}2^{1+a_t}\Gamma(1+a_t)
 \le C e^{-cn/X}.
 \end{split}
 \label{weights:retained-exponential}
\end{equation}
Here \(a_t=(k-t)/(\ell\delta_J)\le\ell^{-3/10}\), so
\(2^{1+a_t}\Gamma(1+a_t)\) is bounded independently of \(k\).
The same Leibniz calculation now proves
\eqref{weights:terminal-derivatives}. All integrands used for
differentiation are dominated on compact sets by the displayed
integrable majorants, which also justifies
\eqref{weights:power-derivative} under the integral.
\end{proof}

\begin{corollary}[Comparison interface]
\label{weights:comparison-interface}
For every fixed derivative order \(r\), the short weights satisfy
the smoothness assumptions of Proposition~\ref{analytic:comparison}
with derivative bounds \(\ell^{C_r}\), uniformly in \(j,k,X\).
They have length \(X^{1-\Delta_j}\), with
\[
 \Delta_j=\delta_j/4,\qquad
 \tfrac14\ell^{-1/10}\le\Delta_j\le\tfrac14 .
\]
The deterministic term
\begin{equation}
 s=k/\ell,\qquad A_x=(C_F|x|)^s\Gamma(1+s)
 \label{weights:Ax}
\end{equation}
has uniformly bounded derivatives of each fixed order on
\(\mathcal K\), and is bounded above and below by positive constants.
The endpoint enlargement at \(n=1\) can be made inside the fixed
range \(n\ge1/2\).
\end{corollary}

\begin{proof}
The bounds \(k\le\ell^{3/5}\) and
\(\delta_j^{-1}\le\ell^{1/10}\) make
\eqref{weights:short-derivatives} polynomial in \(\ell\).
Choose a fixed smooth function \(\chi\) equal to zero on
\((-\infty,-\log2]\) and equal to one on \([0,\infty)\).
For the comparison alone, replace \(w_{0,x}(u)\) on the real line
by
\[
 \widetilde w_{0,x}(u)=\chi(\ell u)w_{0,x}(u).
\]
It agrees with the original weight at every positive integer
\(n=e^{\ell u}\), and its extra derivatives cost only fixed powers
of \(\ell\), by \eqref{weights:main-derivatives}. Its left support
is \(u\ge-\log2/\ell\), exactly the enlargement \(n\ge1/2\).
For \(j\ge1\), the weights already vanish on a neighborhood of
\(u\le0\). After a fixed smooth dyadic localization \(n=Ny\),
the operator \(y\partial_y\) acts as \(\ell^{-1}\partial_u\), so
the same bounds give all the required seminorms in the scaled
variables. The endpoint extension has therefore taken place in
\(\log n\); no extension to a fixed negative value of \(u\) is used.

Finally \(0<s\le\ell^{-2/5}\), while \(C_F|x|\) lies in a fixed
positive compact interval. The formula for \(A_x\), its derivatives
in \(x\), and the continuity and positivity of \(\Gamma(1+s)\)
give the assertions concerning \(A_x\).
\end{proof}

\subsection{The terminal mean square}

We use the following consequence of Proposition~\ref{analytic:polynomial-moment}.
For general $X$, cover $[X,2X)$ by at most two dyadic blocks whose
heights are comparable to $X$; the normalization and polynomial factors
in that proposition then change by fixed constants only.
For any fixed smooth compactly supported function \(G\) on
\((0,\infty)\), there are fixed
exponents \(A,B\) such that
\begin{equation}
 \frac1X\sum_{\substack{d\ {\rm admissible}\\X\le|d|<2X}}
 \left|\sum_{n\ge1}\frac{\lambda(n)\chi_d(n)}{\sqrt n}\,
                         n^{-it}G(n/N)\right|^2
 \ll_G \ell^B(1+|t|)^A(1+N/X)^A.
 \label{weights:moment-interface}
\end{equation}
Its applicability includes the fixed enlarging bump used below.

\begin{proposition}[Terminal mean square]
\label{weights:terminal}
There is a fixed exponent \(C_T\), chosen before \(k_0\), such that
uniformly in \eqref{weights:range},
\begin{equation}
 \frac1X\sum_{\substack{d\ {\rm admissible}\\X\le|d|<2X}}
       |S_{*,\chi_d}(d/X)|^2
       \ll \ell^{C_T}\delta_J^{2k}.
 \label{weights:terminal-square}
\end{equation}
The same estimate holds after restricting the summation to either
value of \(\chi_d(2)\).
\end{proposition}

\begin{proof}
Choose a smooth dyadic partition with a fixed
\(\rho\in C_c^\infty((1/2,2))\) satisfying
\[
 \sum_{N\in\{1,2,4,\ldots\}}\rho(n/N)=1\qquad(n\ge1).
\]
For example one may start with a smooth nonincreasing function
equal to one on \((0,1]\) and zero on \([2,\infty)\), and take
its difference at arguments \(y\) and \(2y\); choosing its
transition strictly inside \((1,2)\) gives the asserted compact
support. Fix also \(G\in C_c^\infty((0,\infty))\) equal to one
on \(\operatorname{supp}\rho\). Put
\[
 T_N(y,x)=\rho(y)w_{*,x}(\ell^{-1}\log(Ny)),\qquad
 \widehat T_N(t,x)=\int_0^\infty T_N(y,x)y^{it}\frac{dy}{y}.
\]
Select once and for all integers \(r>A/2+2\) and \(D>A/2+2\),
and then choose \(k_0\ge r+2\). If derivatives required elsewhere
are of higher order, increase this fixed \(r\) first.
By Lemma~\ref{weights:derivatives}, differentiation on the
\(y\)-scale through order \(r\) gives
\[
 \left|(y\partial_y)^a T_N(y,x)\right|
       \ll_r \ell^{C_r}\delta_J^k e^{-c'N/X}
       \qquad(0\le a\le r),
\]
uniformly in \(x\in\mathcal K\).
There is no boundary term when integrating by parts in \(\log y\).
Consequently,
\begin{equation}
 \sup_{x\in\mathcal K}|\widehat T_N(t,x)|
   \ll_{r,D}\ell^{C_r}\delta_J^k
       (1+N/X)^{-D}(1+|t|)^{-r}.
 \label{weights:terminal-kernel}
\end{equation}
In passing from the exponential to \((1+N/X)^{-D}\), the constant
depends on the fixed \(D\), not on \(k\).

Mellin inversion now expresses the \(N\)-piece of the terminal sum as
\[
 \frac1{2\pi}\int_{\mathbb R}
   \widehat T_N(t,d/X)N^{it}
   P_d(N,t)\,dt,\qquad
 P_d(N,t)=\sum_{n\ge1}\frac{\lambda(n)\chi_d(n)}{\sqrt n}\,
                        n^{-it}G(n/N).
\]
Indeed \(G(n/N)T_N(n/N,x)=T_N(n/N,x)\).
Define the normalized family norm by
\[
 \|B_d\|_{2,X}
   =\left(\frac1X\sum_{\substack{d\ {\rm admissible}\\X\le|d|<2X}}
                      |B_d|^2\right)^{1/2}.
\]
We bound the kernel uniformly before applying the family moment.
For each fixed \(t\),
\[
 \|\widehat T_N(t,d/X)P_d(N,t)\|_{2,X}
   \le \sup_{x\in\mathcal K}|\widehat T_N(t,x)|
                           \|P_d(N,t)\|_{2,X}.
\]
Minkowski's inequality, followed by
\eqref{weights:moment-interface} and
\eqref{weights:terminal-kernel}, therefore yields
\[
 \begin{split}
 \|S_{*,\chi_d}(d/X)\|_{2,X}
 &\ll \ell^{C_r+B/2}\delta_J^k
   \sum_{N=1,2,4,\ldots}(1+N/X)^{-D+A/2}
   \int_{\mathbb R}(1+|t|)^{-r+A/2}\,dt\\
 &\ll \ell^{C_r+B/2+1}\delta_J^k.
 \end{split}
\]
The integral converges by the fixed choice of \(r\).
There are \(O(\ell)\) dyadic terms with \(N\le X\), and those
with \(N>X\) form a convergent geometric tail by the choice of
\(D\). All sums and integrals can first be truncated; the displayed
integrable majorants justify passage to the limit. Squaring proves
\eqref{weights:terminal-square}, for example with
\(C_T=2C_r+B+2\). This exponent depends only on choices made before
\(k_0\). Restricting the family can only decrease the nonnegative
sum.
\end{proof}

\begin{lemma}[Uniform absorption of the terminal loss]
\label{weights:terminal-small}
After increasing the fixed \(k_0\), one sufficiently large lower
bound on \(X\) makes
\begin{equation}
 \frac1X\sum_{\substack{d\ {\rm admissible}\\X\le|d|<2X}}
       |S_{*,\chi_d}(d/X)|^2
       \ll 2^{-3Jk/2}
 \label{weights:terminal-small-square}
\end{equation}
simultaneously for all \(k\) in \eqref{weights:range}.
Moreover, for any fixed \(c,L>0\), the comparison errors obey
\begin{equation}
 e^{-c\Delta_j\ell}
       \le k^{-L}2^{-Ljk}\qquad(0\le j\le J)
 \label{weights:comparison-absorption}
\end{equation}
for one sufficiently large lower bound on \(X\).
\end{lemma}

\begin{proof}
The ratio of the right side of
\eqref{weights:terminal-square} to \(2^{-3Jk/2}\), apart from its
fixed implied constant, is
\[
 \ell^{C_T}2^{-Jk/2}.
\]
For \(J\ge1\) this decreases with \(k\), so throughout the permitted
range it is at most its value at the fixed integer \(k_0\).
The floor in the definition of \(J\) gives
\[
 \ell^{C_T}2^{-Jk/2}
       \le \ell^{C_T}2^{-Jk_0/2}
       \le 2^{k_0/2}\ell^{C_T-k_0/20}.
\]
Choose \(k_0>20C_T\) after \(C_T\) is fixed. The last expression
then tends to zero with \(X\). In particular the choice of a
single \(X\)-threshold works for the entire \(k\)-range.

For the second assertion,
\(\Delta_j\ell\ge\ell^{9/10}/4\), whereas uniformly in the same
range
\[
 L\log k+Ljk\log2
       =O_L(\ell^{3/5}\log\ell)=o(\ell^{9/10}).
\]
Taking logarithms proves \eqref{weights:comparison-absorption}.
\end{proof}
